\todo{Lecture 22}
\begin{reminder}
\begin{lemma}
Soient $x,y\in \mathbf{R}$. 
$$
e^{x+y}=e^{ x }e^{ y }.
$$
\end{lemma}
\begin{proof}[Prueve-brouillon]
\begin{align*}
\left( \sum_{i=0}^{N} \frac{x^{i}}{i!} \right)\left( \sum_{i=0}^{N} \frac{y^{i}}{i!} \right) & =\sum_{k=0}^{2N} \sum _{i+j=k} \frac{x^{i}}{i!} \frac{y^{j}}{j!} \\
 & =\sum_{k=0}^{2N} \sum_{j=0}^{k} \binom{ k}{j } \frac{x^{k-j}y^{j}}{k!} \\
 & =\sum_{k=0}^{2N} \frac{(x+y)^{k}}{k!} \\
 & \overset{N\to \infty}=e^{(x+y)}. 
\end{align*}
\end{proof}
\end{reminder}
\begin{lemma}
Soit $A,B\in \mathbf{C}^{n\times n}$ tels que $AB=BA$. Alors
$$
e^{A+B}=e^{ A }e^{ B }.
$$
\end{lemma}
\begin{proof}
exercise.
\end{proof}
\begin{definition}
Une matrice $N\in \mathbf{C}^{n\times n}$ est nilpotente, s'il existe un entier $k$ tel que $N^{k}=0$.
\end{definition}
\begin{theorem}[Forme normale de Jordan]
Soit $A\in \mathbf{C}^{n\times n}$. Ils existent $P$, $\mathrm{diag}(\lambda_{1},\dots,\lambda _{n})$ et $N$ dans $\mathbf{C}^{n\times n}$ tels que $P$ est inversible, $N$ est une matrice nilpotente, $N$ et $\mathrm{diag}(\lambda_{1},\dots,\lambda _{n})$ commutent et 
$$
A=P^{-1}(\mathrm{diag}(\lambda_{1},\dots,\lambda _{n})+N)P.
$$
\end{theorem}

Maintenant, on veut calculer $e^{ tA }$. On a 
\begin{align*}
A^{i} & =(P^{-1}(\mathrm{diag}(\lambda_{1},\dots,\lambda _{n})+N)P)^{i} \\
 & =P^{-1}\underbrace{ (\mathrm{diag}(\lambda_{1},\dots,\lambda _{n})+N)^{i} }_{ J^{i} }P.
\end{align*}
Alors
$$
e^{ A }=P^{-1}e^{ J }P
$$
et 
$$
e^{ tA }=P^{-1}e^{ tJ }P=P^{-1}e^{ t\mathrm{diag}(\lambda_{1},\dots,\lambda _{n}) }e^{tN}P.
$$
On a
$$
e^{t\mathrm{diag}(\lambda_{1},\dots,\lambda _{n})}=\mathrm{diag}(e^{ t\lambda_{1} },\dots,e^{ t\lambda _{n} })
$$
et
$$
e^{ tN }=I + tN +\dots + \frac{t^{k-1}N^{k-1}}{(k-1)!}
$$
ou $N^{k}=0$.
\begin{example}[Continuation de \ref{exe2101}]
$A$ est diagonalisable. On a
$$
P=\begin{pmatrix}
1 & 1  \\
i & -i
\end{pmatrix}, \quad D=\begin{pmatrix}
-i & 0  \\
0 & i
\end{pmatrix}.
$$
Les valeurs propres de $A$, $\lambda_{1}=-i$ et $\lambda_{2}=i$, ont pour vecteurs propres les colonnes $p_{1}$ et $p_{2}$ de $P$. 

On calcule les $N$ premiers termes de $e^{ tA }$ 
\begin{align*}
\sum_{k=0}^{N} t^{k} \frac{A^{k}}{k!} & =P \sum_{k=0}^{N} \begin{pmatrix} \frac{(-it)^{k}}{k!} & 0 \\
0 & \frac{(it)^{k}}{k!}
\end{pmatrix}P^{-1} \\
 & = P\begin{pmatrix}
e^{ -it } & 0 \\
0 & e^{ it }
\end{pmatrix}P^{-1} \\
 & =P\begin{pmatrix}
\cos t-i\sin t & 0 \\
0 & \cos t+i\sin t 
\end{pmatrix}P^{-1} \\
 & =\begin{pmatrix}
\cos t & -\sin t \\
\sin t & \cos t
\end{pmatrix}.
\end{align*}
\end{example}

\chapter{La forme normale de Jordan}
\begin{definition}
Un bloc Jordan est define comme
$$
J(\lambda,n)=\begin{pmatrix}
\lambda & 1 &  &  &  \\
 & \lambda & 1 &  & \\
 &  &  \ddots &  \ddots\\
 &  &  &  \lambda & 1 \\
 &  &  &  & \lambda
\end{pmatrix}\in \mathbf{C}^{n\times n}
$$
ou $\lambda \in \mathbf{C}^{n\times n}$ et $n\in \mathbf{N}$. On a $J(\lambda,n)=\mathrm{diag}(\lambda,\dots,\lambda)+J(0,n)$.
\end{definition}

On a 
$$
J(0,n)\begin{pmatrix}
x_{1} \\
\vdots \\
x_{n}
\end{pmatrix}=\begin{pmatrix}
x_{2}  \\
\vdots  \\
x_{n} \\
0
\end{pmatrix}
$$
alors
$$
J(0,n)^{i}\begin{pmatrix}
x_{1} \\
\vdots  \\
x_{n}
\end{pmatrix}=\begin{pmatrix}
x_{i+1} \\
\vdots \\
x_{n} \\
0 \\
\vdots \\
0
\end{pmatrix}.
$$

La matrice $J(0,n)$ est nilpotente.
\begin{definition}
Une matrice $A\in \mathbf{C}^{n\times n}$ est une forme normale de Jordan si
$$
A=\begin{pmatrix}
J(\lambda_{1},n_{1}) &  &  \\
 &  \ddots &  \\
 &  &   J(\lambda _{k},n_{k} )
\end{pmatrix}
$$
est une matrice a blocs diagonale ou chaque bloc est de Jordan.
\end{definition}
\begin{remark}
$$
f_{A}(x)=(-1)^{n}\prod(x-\lambda _{i})^{n_{k}}.
$$
\end{remark}